%% CHEOPS Workshop @ EuroSys 2026
%% ACM sigplan format (2-column, 5-6 pages)
\documentclass[sigplan,nonacm,10pt]{acmart}

%% Remove ACM-specific metadata for workshop submission
\settopmatter{printacmref=false}
\renewcommand\footnotetextcopyrightpermission[1]{}
\pagestyle{plain}

\usepackage{booktabs}
\usepackage{subcaption}
\usepackage{xspace}
\usepackage{listings}
\usepackage{balance}

%% Listings style for code
\lstset{
  basicstyle=\ttfamily\small,
  breaklines=true,
  frame=single,
  numbers=left,
  numberstyle=\tiny,
  language=Python,
}

%% Shortcuts
\newcommand{\sysname}{CephSem\xspace}
\newcommand{\eg}{e.g.,\xspace}
\newcommand{\ie}{i.e.,\xspace}
\newcommand{\cf}{cf.\xspace}

\begin{document}

\title{\sysname: LLM-Driven Semantic Object Management for Ceph RADOS}

\author{Author Name}
\affiliation{
  \institution{Institution}
  \city{City}
  \country{Country}
}
\email{author@example.com}

\begin{abstract}
Modern distributed storage systems like Ceph manage petabytes of unstructured
data through low-level interfaces (CLI, API) that require administrators to know
exact object names, pool structures, and command syntax.
We present \sysname, a system that integrates large language models (LLMs) with
Ceph's RADOS layer to enable semantic, natural-language object management.

\sysname makes three contributions:
(1)~a \emph{Ceph-native embedding store} that uses RADOS extended attributes
(xattrs) to co-locate semantic embeddings with object data, eliminating external
vector database dependencies;
(2)~a \emph{RADOS watch/notify--based change detector} that leverages Ceph's
native notification mechanism for near-instant index updates instead of polling;
and
(3)~an \emph{intent-classification pipeline} that maps natural-language queries
to RADOS operations with \textbf{XX.X\%~$\pm$~X.X\%} intent accuracy and
\textbf{XXXX~$\pm$~XXX\,ms} end-to-end latency, where LLM inference accounts
for XX\% of the total.

We evaluate \sysname on a 3-node, 29-OSD Ceph cluster with XXX test cases
over N~repeated trials, comparing against CLI baselines and analyzing
scalability from 10 to 1{,}000 objects.
\end{abstract}

\keywords{Ceph, RADOS, LLM, semantic search, object storage}

\maketitle

%% ============================================================
\section{Introduction}\label{sec:intro}
%% ============================================================

% Para 1: Problem
Managing large-scale object stores requires administrators to remember exact
object names, navigate hierarchical namespaces, and compose complex CLI
commands~\cite{weil06ceph}.
As storage clusters grow to millions of objects across hundreds of pools,
discoverability becomes a significant operational burden.

% Para 2: Opportunity
Recent advances in large language models (LLMs) and embedding-based retrieval
make it feasible to layer a natural-language interface on top of existing storage
APIs.
However, naively wrapping a storage system with an LLM chat interface risks
being a mere ``application'' rather than a systems contribution.

% Para 3: Our approach
We argue that the integration must be \emph{Ceph-native}: embeddings should be
stored \emph{inside} RADOS using extended attributes (xattrs), changes should be
detected via RADOS watch/notify rather than polling, and the semantic index
should benefit from Ceph's replication and failure handling.

% Para 4: Contributions
This paper makes three contributions:
\begin{enumerate}
  \item \textbf{RADOS xattr embedding store.}
    We store embedding vectors (384-dimensional, float32) as extended attributes
    on the objects themselves (\S\ref{sec:design:xattr}).
    This co-locates semantic metadata with data, inherits Ceph's replication
    factor, and eliminates an external ChromaDB/Milvus dependency.

  \item \textbf{Watch/notify change detection.}
    We replace poll-based pool scanning with RADOS \texttt{watch2}/\texttt{notify}
    on a sentinel object (\S\ref{sec:design:watch}).
    This reduces change-detection latency from 60\,s to sub-second with
    $O(1)$ I/O per change instead of $O(N)$ per scan.

  \item \textbf{Intent-classification evaluation.}
    We present a systematic evaluation of LLM-based intent classification for
    storage operations, with XXX test cases across XX~categories, N~repeated
    trials, latency decomposition, scalability analysis, and CLI baseline
    comparison (\S\ref{sec:eval}).
\end{enumerate}


%% ============================================================
\section{Background \& Motivation}\label{sec:background}
%% ============================================================

\subsection{Ceph and RADOS}

Ceph is a distributed storage system that provides object, block, and file
storage on a unified platform~\cite{weil06ceph}.
At its core, the Reliable Autonomic Distributed Object Store (RADOS) manages
data placement, replication, and recovery across OSDs.
RADOS supports extended attributes (xattrs) on every object---key-value pairs
stored alongside the object data with the same durability guarantees.
RADOS also provides a \texttt{watch}/\texttt{notify} mechanism: a client can
\emph{watch} an object and receive asynchronous callbacks when another client
\emph{notifies} the same object.

\subsection{LLMs for Systems Management}

Recent work has explored LLMs for database query optimization~\cite{%TODO},
log analysis~\cite{%TODO}, and cloud operations~\cite{%TODO}.
However, these systems typically treat the LLM as an external oracle that
generates commands.
In contrast, we integrate LLM-derived metadata \emph{into} the storage layer
via RADOS xattrs, making the semantic index a first-class citizen of the
storage system.

\subsection{Semantic Search and Embeddings}

Transformer-based sentence encoders (\eg Sentence-BERT~\cite{reimers2019sentencebert})
map text to dense vectors where cosine similarity approximates semantic
relatedness.
For storage management, embeddings enable operators to query objects by meaning
(``find my Kubernetes configs'') rather than by name or prefix.


%% ============================================================
\section{Design}\label{sec:design}
%% ============================================================

Figure~\ref{fig:arch} shows \sysname's architecture.
A natural-language query flows through four stages:
\emph{intent classification} (LLM), \emph{embedding generation}
(Sentence-BERT), \emph{vector search} (RADOS xattr scan or ChromaDB),
and \emph{RADOS operation execution}.

\begin{figure}[t]
  \centering
  %% Replace with actual figure
  \fbox{\parbox{0.9\columnwidth}{\centering\vspace{2em}
    [Architecture diagram placeholder]\\
    Query $\rightarrow$ LLM $\rightarrow$ Embedding $\rightarrow$ RADOS xattr search $\rightarrow$ Result
  \vspace{2em}}}
  \caption{\sysname architecture.  Solid arrows show the query path;
           dashed arrows show the indexing path.}
  \label{fig:arch}
\end{figure}

\subsection{RADOS Xattr Embedding Store}\label{sec:design:xattr}

Each indexed object receives five extended attributes:
\begin{itemize}
  \item \texttt{user.semantic.embedding} --- binary float32 vector (1.5\,KB for 384-d)
  \item \texttt{user.semantic.model} --- embedding model identifier
  \item \texttt{user.semantic.indexed\_at} --- ISO timestamp
  \item \texttt{user.semantic.content\_hash} --- SHA-256 prefix for staleness detection
  \item \texttt{user.semantic.preview} --- text preview (up to 4\,KB)
\end{itemize}

A manifest object (\texttt{\_semantic\_index\_manifest}) tracks all indexed
object names, enabling efficient enumeration without scanning the full pool.

\paragraph{Search.}
To answer a semantic query, \sysname generates an embedding for the query
text, reads \texttt{user.semantic.embedding} from each indexed object via
\texttt{get\_xattr()}, and computes cosine similarity.
This is a linear scan---$O(N)$ in the number of indexed objects---which is
acceptable for administrative workloads (typically $<$50\,K objects per pool).

\paragraph{Comparison with external vector DB.}
Table~\ref{tab:xattr-vs-chromadb} compares the two approaches.

\begin{table}[t]
\centering
\caption{RADOS xattr store vs.\ external vector database.}
\label{tab:xattr-vs-chromadb}
\begin{tabular}{lcc}
\toprule
\textbf{Property} & \textbf{RADOS Xattr} & \textbf{ChromaDB} \\
\midrule
Deployment        & None (built-in) & Separate process \\
Replication       & Ceph-native     & Must configure \\
Lifecycle         & With object     & Manual cleanup \\
ANN indexing      & No (linear)     & HNSW \\
Search at 1K obj  & XX\,ms          & XX\,ms \\
Search at 10K obj & XX\,ms          & XX\,ms \\
\bottomrule
\end{tabular}
\end{table}


\subsection{RADOS Watch/Notify Change Detection}\label{sec:design:watch}

The original design polled the pool every 60\,s, listing all objects and
comparing modification times---$O(N)$ I/O per cycle.

We replace this with RADOS \texttt{watch2}/\texttt{notify}:
\begin{enumerate}
  \item A sentinel object (\texttt{\_watch\_sentinel}) is created at startup.
  \item The watcher calls \texttt{ioctx.watch()} on the sentinel,
        registering an asynchronous callback.
  \item When any writer creates, updates, or deletes an object, it calls
        \texttt{ioctx.notify()} on the sentinel with a JSON payload
        describing the change.
  \item The callback triggers immediate re-indexing.
  \item A low-frequency fallback poll (every 5\,min) catches changes from
        external writers that do not notify.
\end{enumerate}

This reduces change-detection latency from 60\,s to sub-second and eliminates
per-cycle $O(N)$ scanning.


\subsection{Intent Classification Pipeline}\label{sec:design:intent}

Given a natural-language query, the LLM classifies it into one of
XX~intent types (\eg \texttt{semantic\_search}, \texttt{read\_object},
\texttt{cluster\_health}) and extracts parameters (\eg object names, queries).
The classification prompt is structured as a single LLM completion call
with a JSON schema.

\paragraph{Latency decomposition.}
We instrument the pipeline to measure five phases independently:
LLM inference, embedding generation, vector search, RADOS I/O, and response
formatting.
This decomposition is critical for identifying optimization targets.


%% ============================================================
\section{Implementation}\label{sec:impl}
%% ============================================================

\sysname is implemented in $\sim$5\,K lines of Python.
We use Ollama~\cite{ollama} for local LLM serving (Llama~3.2, 3B parameters),
Sentence-Transformers \texttt{all-MiniLM-L6-v2} (384-d embeddings), and
\texttt{python3-rados} for RADOS access.
The RADOS xattr store (\S\ref{sec:design:xattr}) adds $\sim$300~lines;
the watch/notify watcher (\S\ref{sec:design:watch}) adds $\sim$250~lines.
Both are drop-in replacements for the ChromaDB store and polling watcher
respectively.


%% ============================================================
\section{Evaluation}\label{sec:eval}
%% ============================================================

We evaluate \sysname on a production-like cluster to answer:
\begin{enumerate}
  \item[\textbf{Q1}] How accurately does the LLM classify storage intents?
  \item[\textbf{Q2}] What is the end-to-end latency and where is time spent?
  \item[\textbf{Q3}] How does performance scale with corpus size?
  \item[\textbf{Q4}] What is the overhead vs.\ direct CLI operations?
\end{enumerate}

\subsection{Setup}

\paragraph{Cluster.} 3-node Ceph cluster, 29~OSDs, $\sim$22.9\,TB raw,
Ceph Reef, HEALTH\_OK.
\paragraph{Model.} Llama~3.2 (3B) via Ollama on localhost, GPU: [specify].
\paragraph{Embeddings.} \texttt{all-MiniLM-L6-v2}, 384-d, CPU.
\paragraph{Test suite.} XXX~test cases across XX~categories
(search, read, write, delete, stats, cluster, documentation, complex,
ambiguous, edge cases), three difficulty levels.
\paragraph{Methodology.} Each experiment reports mean~$\pm$~std over
$N=5$ independent runs with fresh agent state.


\subsection{Q1: Intent Classification Accuracy}

% TODO: Fill from multi-run evaluation output
Table~\ref{tab:eval-results} shows the main results.
\sysname achieves \textbf{XX.X\%~$\pm$~X.X\%} intent accuracy and
\textbf{XX.X\%~$\pm$~X.X\%} parameter extraction accuracy.

%% Auto-generated from evaluation results
%% Replace XX placeholders with actual numbers after running benchmarks

% Table: Evaluation Results (mean ± std over N runs)
\begin{table}[t]
\centering
\caption{Agent evaluation results (mean $\pm$ std, N=5 runs, model: llama3.2).}
\label{tab:eval-results}
\begin{tabular}{lr}
\toprule
\textbf{Metric} & \textbf{Result} \\
\midrule
Intent Accuracy (\%) & $XX.X \pm X.X$ \\
Parameter Accuracy (\%) & $XX.X \pm X.X$ \\
Response Quality (\%) & $XX.X \pm X.X$ \\
\midrule
Avg Latency (ms) & $XXXX \pm XXX$ \\
P50 Latency (ms) & $XXXX \pm XXX$ \\
P95 Latency (ms) & $XXXX \pm XXX$ \\
\bottomrule
\end{tabular}
\end{table}



\subsection{Q2: Latency Decomposition}

Table~\ref{tab:latency-decomp} breaks down the end-to-end latency.
LLM inference dominates at XX\%, followed by [second largest].
RADOS I/O accounts for only XX\%, confirming that the overhead is
in the intelligence layer, not the storage layer.


\subsection{Q3: Scalability}

Table~\ref{tab:scalability} shows indexing throughput and search latency
as corpus size increases from 10 to 1{,}000 objects.
Search latency grows [linearly/sublinearly] with the RADOS xattr backend,
as expected from the linear scan.


\subsection{Q4: CLI Baseline Comparison}

Table~\ref{tab:cli-comparison} compares agent latency with direct CLI
commands.
The agent incurs XX\texttimes{} overhead on average, dominated by LLM
inference.
For semantic operations (\eg ``find files about Kubernetes''), the CLI
has no equivalent, representing the unique value of the semantic layer.


%% ============================================================
\section{Discussion}\label{sec:discussion}
%% ============================================================

\paragraph{When to use RADOS xattrs vs.\ external vector DB.}
The xattr backend is ideal for pools with $<$50K objects where simplicity
and data locality are priorities.
For larger corpora requiring approximate nearest neighbor (ANN) search,
an external vector database (\eg ChromaDB, Milvus) is preferable.
\sysname supports both backends and can be configured per-pool.

\paragraph{LLM model choice.}
We use Llama~3.2 (3B) for local inference.
Larger models (\eg 8B, 70B) may improve accuracy but increase latency.
Cloud-hosted models (\eg GPT-4) offer higher accuracy but introduce
network latency and privacy concerns for storage management.

\paragraph{Security considerations.}
The LLM agent operates with the Ceph admin keyring.
Production deployments should use capability-restricted keyrings
and may need to audit LLM-generated operations before execution.

\paragraph{Limitations.}
The linear-scan search in the xattr backend limits scalability.
A Ceph-native ANN index (\eg using RADOS classes) is future work.
The intent classifier depends on prompt engineering and may misclassify
novel query patterns.


%% ============================================================
\section{Related Work}\label{sec:related}
%% ============================================================

\paragraph{LLMs for storage/infrastructure.}
% TODO: Add specific citations
Recent work has applied LLMs to database administration~\cite{%TODO},
cloud resource management~\cite{%TODO}, and log analysis~\cite{%TODO}.
These systems use LLMs as external query planners; \sysname differs by
embedding semantic metadata into the storage layer itself.

\paragraph{Semantic file systems.}
Semantic file systems~\cite{gifford91semantic} organize files by content
attributes rather than directory paths.
\sysname applies this philosophy to object storage, using LLM-generated
embeddings as the semantic index.

\paragraph{Ceph extensions.}
Ceph's RADOS class mechanism allows custom object methods~\cite{weil06ceph}.
Prior work has used RADOS classes for serverless computing~\cite{%TODO}
and data processing~\cite{%TODO}.
Our xattr-based embedding store is complementary and does not require
custom RADOS classes.


%% ============================================================
\section{Conclusion}\label{sec:conclusion}
%% ============================================================

We presented \sysname, a system that integrates LLM-based semantic
understanding with Ceph RADOS.
By storing embeddings as RADOS xattrs and using watch/notify for change
detection, \sysname makes the semantic index a native part of the storage
system.
Our evaluation on a 3-node cluster demonstrates XX\% intent accuracy
with sub-XXX\,ms RADOS overhead per query.
\sysname is open-source at [URL].


%% ============================================================
%% References
%% ============================================================
\bibliographystyle{ACM-Reference-Format}

\begin{thebibliography}{10}

\bibitem{weil06ceph}
S.~A.~Weil, S.~A.~Brandt, E.~L.~Miller, D.~D.~E.~Long, and C.~Maltzahn,
``Ceph: A scalable, high-performance distributed file system,''
in \emph{Proc.\ OSDI}, 2006.

\bibitem{reimers2019sentencebert}
N.~Reimers and I.~Gurevych,
``Sentence-BERT: Sentence embeddings using Siamese BERT-networks,''
in \emph{Proc.\ EMNLP-IJCNLP}, 2019.

\bibitem{gifford91semantic}
D.~K.~Gifford, P.~Jouvelot, M.~A.~Sheldon, and J.~W.~O'Toole~Jr.,
``Semantic file systems,''
in \emph{Proc.\ SOSP}, 1991.

\bibitem{ollama}
Ollama, ``Run large language models locally,''
\url{https://ollama.ai}, 2024.

\end{thebibliography}

\end{document}
